\chapter{Types, variables, and control flow}
\label{ch:types-control}

\begin{quote}
\emph{``In Julia, types are not classes. They are tags that the compiler uses to generate fast code.''}
\end{quote}

% ============================================================
\section{Primitive types}

Julia has a rich hierarchy of numeric types:

\begin{minted}{julia}
# Integers
x = 42              # Int64 (on 64-bit systems)
y = Int32(42)       # explicitly 32-bit
big = Int128(10)^30 # 128-bit integer
huge = big"999999999999999999999999999999999"  # BigInt (arbitrary precision)

typeof(x)           # Int64
sizeof(x)           # 8 bytes

# Floating-point
a = 3.14            # Float64 (double precision)
b = Float32(3.14)   # single precision
c = BigFloat(π)     # arbitrary precision

# Boolean
flag = true         # Bool (subtype of Integer)
flag + 1            # 2 — true is 1, false is 0

# Characters and strings
ch = 'α'            # Char (Unicode character, 4 bytes)
s = "Hello, Julia"  # String (UTF-8 encoded)
\end{minted}

\begin{juliatip}
Use \texttt{typeof(x)} to check any value's type. Use \texttt{supertypes(Int64)} to see the full type hierarchy: \texttt{Int64 <: Signed <: Integer <: Real <: Number <: Any}.
\end{juliatip}

% ============================================================
\section{Variables and assignment}

Variables in Julia are bindings to values, not typed boxes:

\begin{minted}{julia}
x = 10        # x refers to an Int64
x = "hello"   # now x refers to a String — no error

# Multiple assignment
a, b, c = 1, 2.0, "three"

# Swap without a temporary variable
a, b = b, a

# Constants — the compiler optimizes these
const SPEED_OF_LIGHT = 299_792_458  # underscores for readability
const π_approx = 355 // 113        # Rational{Int64}
\end{minted}

\begin{warningbox}
Redefining a \texttt{const} raises a warning (not an error) if the new value has a different type. Use \texttt{const} for values that truly should not change.
\end{warningbox}

% ============================================================
\section{The type hierarchy}

Julia's type system is a tree rooted at \texttt{Any}:

\begin{minted}{julia}
# Abstract types — cannot be instantiated, used for dispatch
abstract type Shape end

# Concrete types — can be instantiated
struct Circle <: Shape
    radius::Float64
end

struct Rectangle <: Shape
    width::Float64
    height::Float64
end

# Check type relationships
Circle <: Shape       # true
Circle <: Any         # true
Int64 <: Number       # true
String <: Number      # false

# Useful introspection
subtypes(Integer)     # [Bool, Signed, Unsigned]
supertype(Float64)    # AbstractFloat
\end{minted}

% ============================================================
\section{Strings}

Strings in Julia are immutable, UTF-8 encoded sequences:

\begin{minted}{julia}
# String basics
s = "Julia is fast"
length(s)             # 13 (number of characters)
sizeof(s)             # 13 (bytes — same here, but differs for Unicode)
s[1]                  # 'J' (Char)
s[end]                # 't'
s[1:5]                # "Julia" (substring)

# String interpolation
name = "World"
greeting = "Hello, $name!"
computation = "2 + 3 = $(2 + 3)"

# Multi-line strings
poem = """
    Roses are red,
    Julia is fast,
    Python is nice,
    but C won't last.
    """

# Common operations
uppercase("julia")          # "JULIA"
replace("Hello", "l" => "r")  # "Herro"
split("a,b,c,d", ",")      # ["a", "b", "c", "d"]
join(["one", "two", "three"], " + ")  # "one + two + three"
occursin("fast", "Julia is fast")     # true

# Regular expressions
m = match(r"(\d+)\s*kg", "Weight: 82 kg")
m.captures[1]              # "82"
\end{minted}

\begin{juliatip}
Indexing into UTF-8 strings by byte position can fail for multi-byte characters. Use \texttt{eachindex(s)} or \texttt{collect(s)} to iterate safely over characters.
\end{juliatip}

% ============================================================
\section{Control flow: conditionals}

\begin{minted}{julia}
# if-elseif-else
temperature = 38.5

if temperature >= 39.0
    status = "high fever"
elseif temperature >= 37.5
    status = "low-grade fever"
else
    status = "normal"
end
println("Status: $status")

# Ternary operator
status = temperature >= 37.5 ? "fever" : "normal"

# Short-circuit evaluation
x = -5
x > 0 || error("x must be positive")  # throws an error
x > 0 && println("x is positive")     # does nothing (no error)
\end{minted}

% ============================================================
\section{Control flow: loops}

\begin{minted}{julia}
# for loops with ranges
for i in 1:5
    println("Iteration $i")
end

# Iterating over collections
fruits = ["apple", "banana", "cherry"]
for fruit in fruits
    println(fruit)
end

# Enumerate (index + value)
for (i, fruit) in enumerate(fruits)
    println("$i: $fruit")
end

# Nested loops with a single for
for i in 1:3, j in 1:3
    i == j && println("($i, $j) is on the diagonal")
end

# while loop
n = 1
while n <= 1000
    n *= 2
end
println("First power of 2 above 1000: $n")  # 1024

# break and continue
for i in 1:100
    i % 7 != 0 && continue   # skip non-multiples of 7
    i > 50 && break           # stop after 50
    println(i)
end
\end{minted}

% ============================================================
\section{Comprehensions}

Comprehensions are the Julian way to build collections concisely:

\begin{minted}{julia}
# Array comprehension
squares = [x^2 for x in 1:10]
# [1, 4, 9, 16, 25, 36, 49, 64, 81, 100]

# With a filter
even_squares = [x^2 for x in 1:10 if iseven(x)]
# [4, 16, 36, 64, 100]

# 2D comprehension → Matrix
multiplication_table = [i * j for i in 1:5, j in 1:5]

# Generator expression (lazy — no memory allocation)
total = sum(1/k^2 for k in 1:1_000_000)
# ≈ π²/6

# Dictionary comprehension
word_lengths = Dict(w => length(w) for w in ["Julia", "Python", "R"])
# Dict("Julia" => 5, "Python" => 6, "R" => 1)

# Set comprehension
unique_remainders = Set(x % 7 for x in 1:100)
\end{minted}

\begin{performancetip}
Generator expressions (without square brackets) are lazy: \texttt{sum(x\^{}2 for x in 1:10\^{}8)} uses almost no memory because it never allocates the array. Always prefer generators when you only need to aggregate a result.
\end{performancetip}

% ============================================================
\section{Tuples and named tuples}

\begin{minted}{julia}
# Tuples — immutable, ordered collections
point = (3.0, 4.0)
point[1]                # 3.0
x, y = point            # destructuring

# Named tuples — tuples with field names
patient = (name="Alice", age=34, bp=120)
patient.name            # "Alice"
patient.age             # 34

# Useful for returning multiple values from functions
function divrem_custom(a, b)
    return (quotient=a ÷ b, remainder=a % b)
end
result = divrem_custom(17, 5)
result.quotient         # 3
result.remainder        # 2
\end{minted}

% ============================================================
\section{Nothing, Missing, and error handling}

\begin{minted}{julia}
# Nothing — the absence of a value (like None in Python)
function maybe_find(collection, target)
    for item in collection
        item == target && return item
    end
    return nothing
end

result = maybe_find([1, 2, 3], 5)
isnothing(result)       # true

# Missing — represents missing data (like NA in R)
data = [1.0, 2.0, missing, 4.0]
sum(data)               # missing (propagates!)
sum(skipmissing(data))  # 7.0

# Exception handling
try
    x = parse(Int, "not_a_number")
catch e
    println("Error: $(typeof(e)) — $(e.msg)")
finally
    println("This always runs")
end
\end{minted}

\begin{exercise}
\begin{enumerate}
  \item Explore the type hierarchy: starting from \texttt{Float64}, use \texttt{supertype()} repeatedly to climb to \texttt{Any}. How many levels are there?
  \item Write a function \texttt{classify\_bmi(weight, height)} that returns a named tuple \texttt{(bmi=..., category=...)} using WHO BMI categories.
  \item Using a comprehension, generate all Pythagorean triples $(a, b, c)$ with $a < b < c \leq 50$ where $a^2 + b^2 = c^2$.
  \item Create a vector with \texttt{missing} values: \texttt{v = [10, missing, 30, missing, 50]}. Compute the mean of the non-missing values.
  \item Write a program that reads a string of DNA nucleotides (e.g., \texttt{"ATCGGATCCA"}) and returns a dictionary counting each nucleotide.
  \item Use string interpolation and a loop to print a formatted multiplication table for numbers 1 through 10.
\end{enumerate}
\end{exercise}

% ============================================================
\section{Chapter summary}

\begin{itemize}
  \item Julia has a rich numeric type hierarchy: \texttt{Int64}, \texttt{Float64}, \texttt{BigInt}, \texttt{Rational}, etc.
  \item Variables are untyped bindings; values carry the type.
  \item Strings are UTF-8, immutable, and support interpolation with \texttt{\$}.
  \item Control flow uses \texttt{if/elseif/else}, \texttt{for}, \texttt{while}, and short-circuit operators.
  \item Comprehensions and generators build collections concisely and efficiently.
  \item \texttt{missing} propagates through computations; use \texttt{skipmissing()} to handle it.
  \item Named tuples are lightweight, immutable structs for returning structured data.
\end{itemize}
